% =====================================================================
% CSE 341: Concepts of Programming Languages
% Project D5 (Part 2): Cumulative Retrospective & Self-Assessment
% Student: 220104004931 Enes Özkan
% Compile with: tectonic 220104004931_D5_P2.tex
% =====================================================================
\documentclass[11pt,a4paper]{article}

\usepackage[utf8]{inputenc}
\usepackage[T1]{fontenc}
\usepackage{mathpazo}
\usepackage[scaled=0.92]{helvet}
\usepackage[varqu,varl]{zi4}
\usepackage[english]{babel}
\usepackage{microtype}
\usepackage[a4paper,margin=2.25cm,top=2.0cm,bottom=2.15cm]{geometry}

\usepackage[table]{xcolor}
\definecolor{ansblue}{HTML}{1E3A8A}
\definecolor{anscaption}{HTML}{6B7280}
\definecolor{anscodebg}{HTML}{F9FAFB}
\usepackage[colorlinks=true, linkcolor=ansblue, urlcolor=ansblue,
            citecolor=ansblue, pdfborder={0 0 0}]{hyperref}

\usepackage[explicit]{titlesec}
\titleformat{\section}
  {\large\bfseries\sffamily\color{black}}
  {\thesection}{0.65em}{#1}
\titlespacing*{\section}{0pt}{11pt}{5pt}
\setlength{\parskip}{4.5pt plus 1pt minus 1pt}
\setlength{\parindent}{0pt}
\linespread{1.08}
\setcounter{secnumdepth}{0}

\usepackage{enumitem}
\setlist[itemize]{leftmargin=1.1em,itemsep=1pt,topsep=2pt,parsep=0pt}

\usepackage{fancyhdr}
\pagestyle{fancy}
\fancyhf{}
\renewcommand{\headrulewidth}{0pt}
\fancyfoot[L]{\footnotesize\color{anscaption}%
  CSE 341: ANS Retrospective \& Self-Assessment (D5, Part 2 Cumulative) \(\cdot\)
  Student 220104004931 Enes Özkan}
\fancyfoot[R]{\footnotesize\color{anscaption}Page \thepage}

\newcommand{\code}[1]{\texttt{\detokenize{#1}}}
\newcommand{\grade}[1]{\textbf{#1}}

\hypersetup{
  pdftitle  = {ANS Retrospective and Self-Assessment (D5, Part 2 Cumulative)},
  pdfauthor = {220104004931 Enes Özkan},
  pdfsubject= {CSE 341 Project Part 2},
  pdfkeywords={DSL, chemistry, retrospective, self-assessment, ANS}
}

\begin{document}

\begin{center}
  {\footnotesize\color{anscaption}CSE 341: Concepts of Programming Languages\par}
  \vspace{3pt}
  {\huge\bfseries\sffamily ANS: Alchemist's Notation Script\par}
  \vspace{4pt}
  {\large Retrospective \& Self-Assessment (D5), Part 2 Cumulative\par}
  \vspace{6pt}
  {\footnotesize\color{anscaption}
    Student: 220104004931 Enes Özkan \(\cdot\)
    Spring 2026 \(\cdot\)
    Gebze Technical University\par}
\end{center}
\vspace{7pt}

\section{What Is Working}

At this point I trust the foundation more than I did at the start, but I do not want to oversell it. The C
lexer and recursive-descent parser pass the current valid and malformed D3 tests, and the diagnostics include
file position, which makes the errors usable. The AST dump is also helpful because it lets me compare the
implementation against the EBNF instead of guessing from the source code alone. The part I like most is that
the chemistry notation still reads naturally: \code{compound water}, \code{H2O}, and \code{CH4 + 2*O2}
all keep the program close to the problem domain.

\section{What Is Hard}

The hard part is not only writing the grammar, but making the lexer and parser behave exactly like that
grammar. Formula literals and identifiers look similar enough that a small lexical decision can affect later
syntax. Units are another example. \code{890.3 kJ/mol} is natural chemistry notation, but the same slash is
also an arithmetic operator. Implementing this in C made the work more explicit: AST nodes, recovery paths,
and memory ownership all have to be handled directly. That is useful for learning, but it also makes small
design changes expensive.

\section{Plan For Part 2}

For Part 2, the priority is execution without drifting away from the design. I need symbol tables, static
scoping, type checking, and an interpreter that respects the precedence and assignment rules from D1. One
specific design fix is output: Part 1 treats \code{print} as a special statement, but it would fit ANS better as
a normal built-in call, in the same style as other helper functions. If we make that change, D1, D2, and D3
must all be updated together. I also need domain-specific checks for reaction and unit mistakes, real outputs
for the D3 programs, and one type error caught before execution.

\section{Retrospective}

If I were starting again, I would decide the type system earlier, before writing so much syntax and report text. The biggest late change was moving from the Part 1 idea of loose or heterogeneous lists to homogeneous \code{List<T>}. I still think that was the right decision, because \code{list[0].atomic_mass} can be checked statically and ANS stays strongly typed. The problem was timing: making that decision late meant D1, D2, and D3 all had to be reviewed again together. I would also add small expression-statement tests earlier. The parser looked fine when expressions were inside \code{let} declarations or conditions, but a bare statement like \code{1 + 2 * 3;} exposed a separate bug.

The main feature I left out was a richer chemistry arithmetic model. ANS stores quantities such as \code{890.3 kJ/mol}, but it does not implement full unit arithmetic like adding same-unit energy values or converting \code{g} to \code{kg}. That would need a deeper quantity system and more type rules. We kept the subset smaller so the type checker, static scoping, reaction balance checks, interpreter, and D3 tests could actually be finished and explained.

The hardest design decision for me was separating syntax from what is actually writable. The grammar can parse member access, indexing, and function calls as expressions, but that does not mean every expression can appear on the left side of \code{=}. Keeping assignment as a statement and rejecting function-call results and formula atom counts as assignment targets made the language safer, but it also forced me to be more precise in the checker and in D1 Section 4.7.

\section{Self-Assessment}

Using the project rubric, my self-assessment is:

\begin{itemize}
  \item \grade{Correctness: B.} The final D2/D3 behavior is much stronger than Part 1: the checker rejects semantic errors and the interpreter runs the three examples. I still give myself B because some important bugs were found during the final audits, not in our first tests.
  \item \grade{Design justification: A.} Sections 4.4, 4.7, and 4.8 now connect ANS choices to Sebesta concepts: operational semantics, static scoping, operand evaluation, assignment, and trade-offs.
  \item \grade{Use of Sebesta: B.} I can explain the Chapter 5--7 vocabulary in my own words, especially scope, lifetime, type equivalence, precedence, and assignment. The weak point was that type-equivalence wording needed several review passes before it became precise.
  \item \grade{Originality: A.} ANS is still a real chemistry DSL, not a renamed calculator: formula literals, reaction sides, balance checking, quantities, and molar-mass examples all belong to the domain.
  \item \grade{Code quality: B.} The C implementation is understandable and modular, but the checker still represents types with simple strings in some places. A dedicated type structure would be cleaner.
  \item \grade{Exam readiness: B.} I can defend my owned areas: operational semantics, expressions and assignment, static scoping, parser behavior, and list-type decisions. Before calling it an A, I still want one more review pass on writable-location edge cases and reaction-side iteration.
\end{itemize}

\end{document}
